\documentclass[11pt]{article}
\usepackage{act-sci}
\settitle{Science Knowledge Primer}
\setsubtitle{Part 5 of 6 \textperiodcentered{} The outside facts ACT assumes}
\settarget{$\geq$ 21 correct}
\settime{18 min}

\newcommand{\kt}[1]{\par\smallskip{\small\renewcommand{\arraystretch}{1.13}\begin{tabularx}{\linewidth}{@{}>{\bfseries\raggedright\arraybackslash}p{0.25\linewidth} >{\raggedright\arraybackslash}X@{}}#1\end{tabularx}}\par}

\begin{document}
\maketitleblock

\begin{abstract}
ACT Science is mainly a reasoning test, but a few questions per form, usually two to four, expect you to bring a basic fact with you: that pH 3 is acidic, that photosynthesis releases oxygen, that $V=IR$, that weather happens in the troposphere. These questions are quick points for students who know the facts and guesses for students who do not. This primer collects what ACT has assumed, at roughly a first-year high-school level, in five reference sheets: biology, chemistry, physics, Earth and space science, and a math and lab toolbox. Each sheet ends with a short self-check. Read one sheet per day, cover the right-hand column of each table and recite it, then do the check.
\end{abstract}

\scoringline

\tableofcontents
\clearpage

%% =====================================================================
\section{Rule PK-1: Biology essentials}

\begin{rulebox}{PK-1}{Cells, genetics, energy, ecology, evolution}
\kt{%
Cell parts & \emph{Nucleus}: holds DNA. \emph{Mitochondria}: cellular respiration, make most ATP. \emph{Ribosomes}: build proteins. \emph{Chloroplasts} (plants, algae): photosynthesis. \emph{Cell membrane}: controls what enters and leaves. \emph{Cell wall}: rigid outer layer in plants, fungi, bacteria. Prokaryotes (bacteria) have no nucleus.\\
DNA $\to$ protein & DNA bases pair A--T and C--G (RNA uses U instead of T). A gene is copied to RNA (transcription), then read by ribosomes to build a protein (translation).\\
Cell division & \emph{Mitosis}: one cell $\to$ two identical cells (growth, repair). \emph{Meiosis}: $\to$ four cells with half the chromosomes (eggs, sperm).\\
Genetics & Dominant allele (capital, $A$) masks recessive ($a$). Homozygous $AA$ or $aa$; heterozygous $Aa$. $Aa\times Aa \to$ 1 $AA$ : 2 $Aa$ : 1 $aa$, so a 3:1 dominant:recessive phenotype ratio. $Aa \times aa \to$ 1:1.\\
Photosynthesis & $6\,\ce{CO2} + 6\,\ce{H2O} + \text{light} \to \ce{C6H12O6} + 6\,\ce{O2}$. Uses CO$_2$, releases O$_2$.\\
Respiration & $\ce{C6H12O6} + 6\,\ce{O2} \to 6\,\ce{CO2} + 6\,\ce{H2O} + \text{energy (ATP)}$. Happens in plants \emph{and} animals.\\
Transport & \emph{Diffusion}: particles move from high to low concentration. \emph{Osmosis}: water moves across a membrane toward the side with \emph{more} dissolved solute. A cell in salty water loses water and shrinks.\\
Enzymes & Proteins that speed up reactions without being used up. Each has an optimum temperature and pH; too hot or too acidic and it \emph{denatures} (loses shape and function).\\
Ecology & Producers (plants) $\to$ primary consumers $\to$ secondary consumers; decomposers recycle nutrients. Only about 10\% of energy passes to each next level. Predator numbers rise and fall a step behind prey numbers.\\
Evolution & Natural selection: individuals vary, variation is inherited, and individuals with helpful traits survive and reproduce more, so those traits become more common over generations.\\
}
\end{rulebox}

\begin{q}{3}{PK-1}{Mitochondria carry out cellular respiration, which produces most of a cell's ATP.}
In which organelle does cellular respiration produce most of a cell's ATP?
\choicesgrid{Nucleus}{Ribosome}{Mitochondrion}{Chloroplast}
\end{q}

\begin{q}{2}{PK-1}{$Aa\times Aa$ gives 1 $AA$ : 2 $Aa$ : 1 $aa$. Only $aa$ shows the recessive trait: 1/4.}
Two plants heterozygous for a trait ($Aa \times Aa$) are crossed. What fraction of offspring is expected to show the recessive phenotype?
\choicesrow{0}{1/4}{1/2}{3/4}
\end{q}

\begin{q}{4}{PK-1}{Photosynthesis turns CO$_2$ and water into glucose and releases oxygen.}
The products of photosynthesis are:
\choicesgrid{oxygen and water.}{carbon dioxide and water.}{glucose and carbon dioxide.}{glucose and oxygen.}
\end{q}

\begin{q}{2}{PK-1}{About 10\% per level: 10{,}000 $\to$ 1{,}000 (rabbit) $\to$ 100 kJ (fox).}
In the food chain grass $\to$ rabbit $\to$ fox, the grass stores \SI{10000}{\kilo\joule} of energy. About how much reaches the fox?
\choicesrow{\SI{10}{\kilo\joule}}{\SI{100}{\kilo\joule}}{\SI{1000}{\kilo\joule}}{\SI{5000}{\kilo\joule}}
\end{q}

\begin{q}{1}{PK-1}{Water moves by osmosis toward the higher solute concentration (outside), so the cell loses water and shrinks.}
An animal cell is placed in water that is much saltier than the cell's interior. The cell will most likely:
\choices{lose water and shrink.}{gain water and swell.}{gain salt and swell.}{stay the same size.}
\end{q}

%% =====================================================================
\section{Rule PK-2: Chemistry essentials}

\begin{rulebox}{PK-2}{Atoms, matter, solutions, gases, reactions}
\kt{%
Atom & Protons ($+$) and neutrons (0) in the nucleus; electrons ($-$) outside. \emph{Atomic number} $=$ number of protons (identifies the element). \emph{Mass number} $=$ protons $+$ neutrons. Isotopes differ in neutrons; ions differ in electrons.\\
Periodic table & Columns (groups) share properties. Metals on the left, nonmetals on the right. Group 1 (alkali metals) very reactive; Group 18 (noble gases) unreactive.\\
Bonds & \emph{Ionic}: metal $+$ nonmetal, electrons transferred (NaCl). \emph{Covalent}: nonmetals share electrons (H$_2$O, CO$_2$).\\
Phase changes & Melting (s$\to$l), freezing (l$\to$s), vaporization/boiling (l$\to$g), condensation (g$\to$l), sublimation (s$\to$g), deposition (g$\to$s). Temperature stays constant \emph{during} a phase change.\\
Density & $\rho = m/V$. Water $\approx$ \SI{1.0}{\gram\per\milli\litre}. Less dense than the fluid $\to$ floats.\\
pH & Below 7 acidic, 7 neutral, above 7 basic. Each step of 1 is a factor of 10 in acidity.\\
Solutions & Solute dissolves in solvent. Molarity $=$ moles of solute $\div$ liters of solution. Most solids dissolve better when hot; gases dissolve \emph{less} when hot.\\
Gas laws & $T$ in kelvin: $K = {}^\circ\mathrm{C} + 273$. Constant $T$: $PV$ constant (halve $V$, double $P$). Constant $P$: $V \propto T$. Constant $V$: $P \propto T$.\\
Reactions & Mass is conserved; balanced equations have equal atoms on each side. \emph{Exothermic} releases heat (surroundings warm); \emph{endothermic} absorbs heat (surroundings cool). Rates increase with temperature, concentration, surface area, and catalysts.\\
Formulas & H$_2$O water, CO$_2$ carbon dioxide, O$_2$ oxygen, N$_2$ nitrogen, NaCl table salt, CH$_4$ methane, C$_6$H$_{12}$O$_6$ glucose, HCl hydrochloric acid, NaOH sodium hydroxide (a base).\\
}
\end{rulebox}

\begin{q}{1}{PK-2}{pH below 7 is acidic.}
A solution with a pH of 3 is:
\choicesrow{acidic.}{neutral.}{basic.}{a pure element.}
\end{q}

\begin{q}{3}{PK-2}{$\rho = 40 \div 8 = \SI{5}{\gram\per\centi\metre\cubed}$, greater than water's 1.0, so it sinks.}
An object has a mass of \SI{40}{\gram} and a volume of \SI{8}{\centi\metre\cubed}. Its density and its behaviour in water are:
\choices{\SI{0.2}{\gram\per\centi\metre\cubed}; it floats.}{\SI{5}{\gram\per\centi\metre\cubed}; it floats.}{\SI{5}{\gram\per\centi\metre\cubed}; it sinks.}{\SI{320}{\gram\per\centi\metre\cubed}; it sinks.}
\end{q}

\begin{q}{3}{PK-2}{At constant temperature $PV$ is constant, so halving $V$ doubles $P$.}
A gas is held at constant temperature while its volume is cut in half. Its pressure is:
\choicesrow{quartered.}{halved.}{doubled.}{quadrupled.}
\end{q}

\begin{q}{2}{PK-2}{Gas $\to$ liquid is condensation.}
The change of a substance from a gas to a liquid is called:
\choicesgrid{evaporation.}{condensation.}{sublimation.}{melting.}
\end{q}

\begin{q}{4}{PK-2}{The atomic number is the number of protons.}
The atomic number of an element equals the number of:
\choicesgrid{neutrons in its nucleus.}{protons plus neutrons.}{electron shells.}{protons in its nucleus.}
\end{q}

%% =====================================================================
\section{Rule PK-3: Physics essentials}

\begin{rulebox}{PK-3}{Motion, force, energy, waves, electricity, heat}
\kt{%
Motion & Speed $= d/t$. Velocity is speed with direction. Acceleration $=\Delta v/t$. Near Earth, falling objects accelerate at $g\approx\SI{9.8}{\metre\per\second\squared}$ (ignoring air resistance).\\
Newton's laws & (1) Objects keep their motion unless a net force acts. (2) $F = ma$. (3) Every force has an equal, opposite partner. Weight $= mg$; mass does not change with location.\\
Energy & Kinetic $KE=\tfrac12 mv^2$ (double the speed $\to$ 4$\times$ KE). Gravitational $PE = mgh$. Total energy is conserved: a falling ball loses PE and gains KE. Work $= Fd$; power $=$ work $\div$ time. Units: joule (J), watt (W).\\
Momentum & $p = mv$; conserved in collisions.\\
Waves & $v = f\lambda$. Frequency $\leftrightarrow$ pitch (sound) or colour (light); amplitude $\leftrightarrow$ loudness or brightness. Sound needs a medium and travels fastest in solids. Light travels at \SI{3e8}{\metre\per\second} in a vacuum.\\
EM spectrum & Radio, microwave, infrared, visible, ultraviolet, X-ray, gamma: frequency and energy increase, wavelength decreases.\\
Circuits & $V = IR$; power $P = IV$. \emph{Series}: one path, same current, resistances add. \emph{Parallel}: several paths, same voltage, total resistance decreases.\\
Heat & Flows from hot to cold by conduction (contact), convection (moving fluid), and radiation. Metals conduct well; wood, air, and plastic insulate.\\
Pressure & $P = F/A$. Pressure in a fluid increases with depth.\\
}
\end{rulebox}

\begin{q}{3}{PK-3}{$a = \Delta v/t = 20/4 = \SI{5}{\metre\per\second\squared}$.}
A car speeds up from rest to \SI{20}{\metre\per\second} in \SI{4}{\second}. Its average acceleration is:
\choicesrow{\SI{0.2}{\metre\per\second\squared}}{\SI{4}{\metre\per\second\squared}}{\SI{5}{\metre\per\second\squared}}{\SI{80}{\metre\per\second\squared}}
\end{q}

\begin{q}{1}{PK-3}{$a = F/m = 6/2 = \SI{3}{\metre\per\second\squared}$.}
A net force of \SI{6}{\newton} acts on a \SI{2}{\kilo\gram} cart. Its acceleration is:
\choicesrow{\SI{3}{\metre\per\second\squared}}{\SI{4}{\metre\per\second\squared}}{\SI{8}{\metre\per\second\squared}}{\SI{12}{\metre\per\second\squared}}
\end{q}

\begin{q}{4}{PK-3}{$v = f\lambda = 5 \times 2 = \SI{10}{\metre\per\second}$.}
A wave has a frequency of \SI{5}{\hertz} and a wavelength of \SI{2}{\metre}. Its speed is:
\choicesrow{\SI{0.4}{\metre\per\second}}{\SI{2.5}{\metre\per\second}}{\SI{7}{\metre\per\second}}{\SI{10}{\metre\per\second}}
\end{q}

\begin{q}{2}{PK-3}{$I = V/R = 12/4 = \SI{3}{\ampere}$.}
A \SI{12}{\volt} battery is connected across a \SI{4}{\ohm} resistor. The current is:
\choicesrow{\SI{0.33}{\ampere}}{\SI{3}{\ampere}}{\SI{16}{\ampere}}{\SI{48}{\ampere}}
\end{q}

\begin{q}{3}{PK-3}{As the ball falls, height drops (PE decreases) and speed rises (KE increases); total energy is conserved.}
A ball is dropped from a rooftop. Ignoring air resistance, as it falls its:
\choices{potential energy and kinetic energy both increase.}{potential energy increases and kinetic energy decreases.}{potential energy decreases and kinetic energy increases.}{potential energy and kinetic energy both decrease.}
\end{q}

%% =====================================================================
\section{Rule PK-4: Earth and space essentials}

\begin{rulebox}{PK-4}{Geology, atmosphere, water, astronomy}
\kt{%
Earth's layers & Crust (thin, solid) $\to$ mantle (hot, slowly flowing rock) $\to$ outer core (\emph{liquid} iron--nickel) $\to$ inner core (solid).\\
Plate tectonics & \emph{Divergent} boundaries spread apart (mid-ocean ridges). \emph{Convergent} collide (mountains, subduction, volcanoes). \emph{Transform} slide past (earthquakes).\\
Rock cycle & \emph{Igneous}: cooled magma or lava. \emph{Sedimentary}: compacted sediment, often layered, holds fossils. \emph{Metamorphic}: changed by heat and pressure. In undisturbed layers, deeper layers are older (superposition).\\
Atmosphere & Mostly N$_2$ (78\%) and O$_2$ (21\%). \emph{Troposphere}: lowest layer, all weather, temperature falls with altitude. \emph{Stratosphere}: ozone layer, temperature rises with altitude. Then mesosphere and thermosphere. Air pressure falls with altitude.\\
Weather & Warm air rises and cool air sinks. Rising air cools, water vapour condenses, clouds form. Warm air holds more water vapour.\\
Water cycle & Evaporation, transpiration (from plants), condensation, precipitation, runoff, infiltration into groundwater.\\
Astronomy & Day $=$ one rotation ($\approx$ 24 h). Year $=$ one revolution around the Sun. \emph{Seasons} come from the \ang{23.5} tilt of Earth's axis, not from distance to the Sun. Moon phases come from how much of the Moon's sunlit half we see ($\approx$ 29.5-day cycle). \emph{Tides} come mainly from the Moon's gravity (about two highs per day).\\
Solar system & Mercury, Venus, Earth, Mars (rocky); Jupiter, Saturn, Uranus, Neptune (giant). The Sun is a star; a light-year is a \emph{distance}.\\
}
\end{rulebox}

\begin{q}{2}{PK-4}{The tilt of Earth's axis changes how directly sunlight strikes each hemisphere through the year.}
Earth has seasons mainly because:
\choices{Earth's distance from the Sun changes greatly during the year.}{Earth's axis is tilted relative to its orbit.}{the Moon blocks sunlight during winter.}{the Sun's energy output changes during the year.}
\end{q}

\begin{q}{3}{PK-4}{Igneous rock forms when magma or lava cools and hardens.}
Rock formed when magma cools and hardens is classified as:
\choicesgrid{sedimentary.}{metamorphic.}{igneous.}{fossiliferous.}
\end{q}

\begin{q}{1}{PK-4}{Nearly all weather occurs in the troposphere, the lowest layer.}
In which layer of the atmosphere does nearly all weather occur?
\choicesgrid{Troposphere}{Stratosphere}{Mesosphere}{Thermosphere}
\end{q}

\begin{q}{4}{PK-4}{Ocean tides are caused mainly by the Moon's gravitational pull.}
Ocean tides on Earth are caused mainly by:
\choices{wind blowing across the ocean surface.}{Earth's rotation alone.}{differences in water temperature.}{the gravitational pull of the Moon.}
\end{q}

\begin{q}{3}{PK-4}{The outer core is liquid; the inner core is solid because of extreme pressure.}
Which layer of Earth is liquid?
\choicesrow{Crust}{Mantle}{Outer core}{Inner core}
\end{q}

%% =====================================================================
\section{Rule PK-5: Math, units and lab toolbox}

\begin{rulebox}{PK-5}{The tools behind the arithmetic}
\kt{%
SI units & length metre (m), mass kilogram (kg), time second (s), temperature kelvin (K) or \si{\celsius}, amount mole (mol), force newton (N), energy joule (J), power watt (W), pressure pascal (Pa), current ampere (A), voltage volt (V), resistance ohm ($\Omega$), frequency hertz (Hz).\\
Prefixes & giga $10^{9}$, mega $10^{6}$, kilo $10^{3}$, centi $10^{-2}$, milli $10^{-3}$, micro $10^{-6}$, nano $10^{-9}$.\\
Conversions & Multiply by a fraction equal to 1: $\SI{2.5}{\kilo\metre}\times\frac{\SI{1000}{\metre}}{\SI{1}{\kilo\metre}} = \SI{2500}{\metre}$.\\
Scientific notation & $a\times10^{n}$ with $1\le a<10$. Small numbers have negative exponents: $0.00045 = 4.5\times10^{-4}$.\\
Proportions & \emph{Direct}: $y = kx$ (straight line through the origin; double $x$, double $y$). \emph{Inverse}: $xy = k$ (triple $x$, $y$ becomes one third). \emph{Linear}: $y = mx + b$ with slope $m = \Delta y/\Delta x$.\\
Percent & Percent change $= (\text{new}-\text{old})/\text{old}\times100\%$. Percent error $= |\text{measured}-\text{accepted}|/\text{accepted}\times100\%$.\\
Lab tools & Graduated cylinder (volume, fairly precise), beaker and flask (holding and mixing, rough volume only), pipette or buret (very precise volume), balance (mass), thermometer or probe, stopwatch, pH meter or indicator, spectrophotometer (light absorbed by a solution), ammeter (current, in series), voltmeter (voltage, in parallel).\\
}
\end{rulebox}

\begin{q}{4}{PK-5}{\SI{1}{\kilo\metre} $=$ \SI{1000}{\metre}, so $2.5\times1000 = \SI{2500}{\metre}$.}
\SI{2.5}{\kilo\metre} is equal to:
\choicesrow{\SI{0.0025}{\metre}}{\SI{25}{\metre}}{\SI{250}{\metre}}{\SI{2500}{\metre}}
\end{q}

\begin{q}{2}{PK-5}{Move the decimal 4 places right to get 4.5, so the exponent is $-4$.}
The number 0.00045 written in scientific notation is:
\choicesgrid{$4.5\times10^{-5}$}{$4.5\times10^{-4}$}{$4.5\times10^{-3}$}{$4.5\times10^{4}$}
\end{q}

\begin{q}{2}{PK-5}{Inverse: $xy$ is constant, so tripling $x$ makes $y$ one third as large.}
If $y$ is inversely proportional to $x$ and $x$ triples, then $y$:
\choicesgrid{becomes one ninth as large.}{becomes one third as large.}{triples.}{becomes nine times as large.}
\end{q}

\begin{q}{1}{PK-5}{Of these, only the graduated cylinder is marked for precise volume. Beakers and flasks give rough estimates.}
Which piece of equipment would measure \SI{15.0}{\milli\litre} of water most precisely?
\choicesgrid{A \SI{25}{\milli\litre} graduated cylinder}{A \SI{250}{\milli\litre} beaker}{A \SI{500}{\milli\litre} Erlenmeyer flask}{A test tube}
\end{q}

\makekey
\end{document}
